\section{Candidate components (open weights, September 2026)}
\label{sec:components}

This section surveys the open, self-hostable candidates for each stage. For every candidate we
record quality evidence, the property that decides its latency, and whether it runs on the laptop
CPU or needs a free GPU. Section~\ref{sec:measure} gives the measurements.
WAXAL~\cite{waxal} (Google, CC-BY-4.0, released February 2026) is now the common Amharic
benchmark: about 190\,h of transcribed, image-prompted \emph{spontaneous} Amharic speech (566
training speakers), a 16\,h test set, and a TTS subset. Most numbers below refer to its test set.

\subsection{Speech recognition (ASR)}
The deciding latency property is the \textbf{decoder type}. \emph{Autoregressive} (AR) decoders
(Whisper, LLM-based ASR) emit one token per step, and Amharic needs many tokens per word. CTC decoders
emit all characters in one parallel pass, so their cost scales with audio length only. The second
property is whether the encoder can run \textbf{incrementally} while the user speaks.

\begin{table}[H]
\centering\footnotesize\setlength{\tabcolsep}{4pt}
\begin{tabularx}{\textwidth}{@{}L{3.5cm}rL{1.2cm}cY@{}}
\toprule
\textbf{Model} & \textbf{Params} & \textbf{Decoder} & \textbf{WAXAL WER/CER} & \textbf{Latency-relevant notes} \\
\midrule
Whisper large-v3, zero-shot & 1.55B & AR & 254 / 145\lit & Unusable without fine-tuning~\cite{waxalnet}. \\
MMS-1B, zero-shot & 1B & CTC & 101 / 85\lit & Unusable zero-shot~\cite{waxalnet}. \\
Omnilingual ASR 1B, zero-shot & 1B & CTC & 51.1 / 17.4\lit & 1{,}600+ languages; audio $\le$40\,s; GPU RTF 0.001--0.006 for CTC~\cite{omniasr,waxalnet}. \\
\midrule
\hfl{waxal-benchmarking/whisper-tiny-waxal-amh}{Whisper-tiny-WAXAL} & 39M & AR & 41.3 / 16.8\lit & Fast on CPU; fixed 30\,s encoder window. \\
\hfl{waxal-benchmarking/whisper-small-waxal-amh}{Whisper-small-WAXAL} & 244M & AR & 33.6 / 12.9\lit & Best Whisper option; 3.8\,s encoder floor on our CPU\meas. \\
\hfl{waxal-benchmarking/mms-300m-waxal-amh}{MMS-300M-WAXAL} & 300M & CTC & 37.8 / 13.2\lit & Cost $\propto$ audio length; chunked streaming possible. \\
\hfl{waxal-benchmarking/omniasr-ctc-300m-waxal-amh}{OmniASR-CTC-300M-WAXAL} & 300M & CTC & n/a & fairseq2 checkpoint; same family as Omnilingual. \\
Ethio-ASR (AfriHuBERT) & 94M & CTC & 30.95 / --\lit & Small and accurate; checkpoint availability to confirm~\cite{ethioasr}. \\
Ethio-ASR (MMS-1B) & 1B & CTC & 26.14 / --\lit & Larger multilingual variant~\cite{ethioasr}. \\
\hfl{badrex/Ethio-ASR-amharic}{Ethio-ASR-amharic} (w2v-BERT 2.0) & 600M & CTC & \textbf{22.37} / --\lit & Most accurate open Amharic ASR found; 2.4\,GB fp32; T4-friendly. \\
\hfl{b1n1yam/qwen3-asr-0.6b-amharic-mixed}{Qwen3-ASR-0.6B-Amharic} & 0.6B & AR (LLM) & 36.6 / 15.6\lit & Native 2\,s-chunk streaming and 92\,ms TTFT on GPU~\cite{qwen3asr}; the Qwen tokenizer costs 6.5 tokens/word\meas. \\
Whisper-large-v3-turbo Amharic (Bell) & 809M & AR (4 layers) & n/a & Gated. Cheap turbo decoder, but the large encoder is expensive on CPU. \\
Gemma~4 E2B/E4B audio encoder & $\sim$300M enc. & AR (LLM) & unknown & Half-cascade (Section~\ref{sec:how}); Amharic audio untested. \\
\bottomrule
\end{tabularx}
\caption{ASR candidates (model names link to Hugging Face). Most WAXAL numbers are from the WAXAL-NET
benchmark~\cite{waxalnet} or the model cards; WER normalisation differs between sources. Our own
numbers (Section~\ref{sec:m-asr}) use one normaliser for all models.}
\label{tab:asr-cands}
\end{table}

\paragraph{Runtimes for ASR.} CTranslate2 / faster-whisper (int8 on CPU, fp16 on GPU),
whisper.cpp (supports \code{--audio-ctx} to shrink the fixed 30\,s encoder window for short
utterances), ONNX Runtime and OpenVINO for wav2vec2-family CTC models (OpenVINO targets exactly this
kind of Intel CPU), sherpa-onnx for streaming CTC/transducer models with a built-in VAD, and vLLM for
Qwen3-ASR on GPU.

\subsection{Response generation (LLM)}
Two properties dominate latency: \textbf{time to first token} (prefill, which scales with prompt length and
model size and is largely removed by prompt caching) and \textbf{tokens per Amharic word} (decode
steps). A model that knows Amharic but emits 7 tokens per word is, on the same hardware, about twice as
slow per spoken word as one that emits 3.3.

\begin{table}[H]
\centering\footnotesize\setlength{\tabcolsep}{4pt}
\begin{tabularx}{\textwidth}{@{}L{3.6cm}L{1.7cm}cY@{}}
\toprule
\textbf{Model} & \textbf{Params} & \textbf{Tok/word}\meas & \textbf{Amharic evidence and latency notes} \\
\midrule
\hfl{rasyosef/Llama-3.2-1B-Amharic-Instruct}{Llama-3.2-1B-Amharic-Instruct} & 1.24B & \textbf{1.70} & Continued pre-training on Amharic plus an extended vocabulary. Fluent but often factually wrong in our test\meas. 14 tok/s decode on the laptop\meas. \\
Llama-3.2-400M / 180M-Amharic-Instruct & 0.4 / 0.18B & -- & Weaker. Candidate \emph{draft} models for speculative decoding if their vocabulary matches the 1B model. \\
\hfl{google/gemma-4-E2B-it}{Gemma~4 E2B-it} & 2.3B eff.\ (5.1B) & 3.26 & Pre-trained on 140+ languages~\cite{gemma4card}. Q4 GGUF 3.1\,GB. Ships an MTP draft head for speculative decoding. Amharic quality to be tested. \\
Gemma~4 E4B-it & 4.5B eff.\ (8B) & 3.26 & Stronger; realistic on a T4 in 4-bit. \\
Gemma~3 4B / 12B-it & 4 / 12B & 3.26 & Gemma~2 27B was the best open model on most AfroBench languages~\cite{afrobench}; 12B fits a T4 in 4-bit. \\
Qwen3 / Qwen3.5 & 0.6--4B & 6.48 / 7.08 & Good general ability, but about 2$\times$ the tokens of Gemma per Amharic word. \\
Qwen3.5-Amharic-4B-thinking (community) & 4B & 7.08 & Amharic fine-tune. ``Thinking'' must be off: hidden reasoning tokens delay the first spoken word. \\
Llama~3.2 (base) & 1--3B & 12.18 & Byte-level for Amharic; slow and weak. \\
Phi-4-mini & 3.8B & 9.14 & Not suitable. \\
\bottomrule
\end{tabularx}
\caption{LLM candidates. The tokens-per-word column is measured on 3{,}000 Amharic news sentences
(Section~\ref{sec:m-tok}).}
\label{tab:llm-cands}
\end{table}

\paragraph{Translation cascade (considered and rejected).} Translating Amharic$\rightarrow$English,
answering in English, and translating back (e.g.\ with NLLB) adds two sequence-to-sequence passes
before the first Amharic word exists. That adds about 2$\times$ the latency, and the rigid translations
hurt naturalness. We keep it only as a quality fallback, not as a latency candidate.

\paragraph{Runtimes for LLMs.} llama.cpp (GGUF; CPU and CUDA; prompt cache; speculative decoding
including Gemma~4 MTP heads; grammar-constrained output), vLLM and SGLang (GPU; prefix caching;
continuous batching for concurrency; fp16 on a T4, which has no bf16), and TensorRT-LLM.

\subsection{Speech synthesis (TTS)}
The deciding property is \textbf{time to first playable audio} for the first chunk. Non-autoregressive
models (VITS, diffusion) must synthesise a whole chunk before playing, so a \emph{short first chunk}
is the key lever. Codec-LM models (Orpheus + SNAC) can stream audio frame by frame, but they need
roughly 80+ codec tokens per second of audio, which means a GPU.

\begin{table}[H]
\centering\footnotesize\setlength{\tabcolsep}{4pt}
\begin{tabularx}{\textwidth}{@{}L{3.3cm}rL{2.8cm}Y@{}}
\toprule
\textbf{Model} & \textbf{Params} & \textbf{Architecture} & \textbf{Notes} \\
\midrule
\hfl{facebook/mms-tts-amh}{MMS-TTS-amh} & 36M & VITS (non-autoregressive) & Input romanised with \code{uroman}; CC-BY-NC-4.0 (non-commercial). RTF 0.57 on the laptop\meas. Drops digits and mispronounces English\meas. \\
\hfl{african-low-resource/omnivoice-amharic}{OmniVoice-Amharic} (Gheero-Leyu fork: \hfl{gheero-Leyu/amharic-omnivoice-tts}{link}) & 613M & Discrete diffusion (Qwen3-0.6B + HiggsAudioV2 codec) & Trained on $\sim$331\,h of Amharic; zero-shot voice cloning from 10\,s; Apache-2.0; about 3\,GB VRAM on a T4. The base OmniVoice reaches RTF 0.025 on an H100 with 32 steps, and 16 steps is faster~\cite{omnivoice}. MOS and CER are still ``TBD'' on the card. \\
\hfl{Bmancman/amharic-get-tts}{Orpheus-3B-Amharic} (+ GGUF) & 3B & LLM $\rightarrow$ SNAC codec tokens & True frame-level streaming; needs a strong GPU for real time; 6.6\,GB. \\
\hfl{AddisuSeteye/speecht5_tts_amharic}{SpeechT5-Amharic} & 144M & SpeechT5 + HiFi-GAN & Trained on ALFFA; autoregressive spectrogram decoder. \\
Piper voice trained on WAXAL-TTS & $\sim$20M & VITS (ONNX) & No Amharic voice exists yet. Piper is built for fast CPU inference, so this is a strong candidate for the CPU path (E-TTS-6). \\
\bottomrule
\end{tabularx}
\caption{TTS candidates. The Gheero-Leyu fork of OmniVoice-Amharic suggests it may already be the
organisation's preferred voice. Its speed on a T4 is the most important unknown for Week~7.}
\label{tab:tts-cands}
\end{table}

\subsection{VAD and turn detection}
\begin{itemize}
  \item \textbf{Silero VAD} (MIT; 2\,MB ONNX): 0.22\,ms per 32\,ms frame on the laptop\meas. Negligible cost.
  \item \textbf{TEN VAD}: a newer low-latency VAD that claims faster speech-to-silence transitions; worth a direct comparison.
  \item \textbf{Smart Turn v3.2} (BSD-2; 8M params; Whisper-tiny encoder): audio-only semantic end-of-turn
  for 14 languages. 120\,ms per call on our CPU\meas; \textbf{does not work on Amharic}\meas\ (Section~\ref{sec:m-endpoint}).
  Its training code is open, so an Amharic fine-tune on WAXAL is feasible (E-EP-3).
  \item \textbf{Text-based end-of-utterance models}: a small LM scores whether the partial transcript
  is a complete sentence. They need an Amharic model, which could be the 180M Amharic Llama.
\end{itemize}

\subsection{Orchestration and transport}
\textbf{Pipecat} (Python; frame-based pipelines; built-in interruption handling, Smart Turn and Silero
integration) and \textbf{LiveKit Agents} (WebRTC) are the open frameworks used in production voice
agents. A plain \code{asyncio} WebSocket server with 20\,ms 16\,kHz PCM frames is simpler and makes
every millisecond visible. For a single-machine demo, in-process queues avoid all network cost. Our
recommendation is to build the measurement harness on plain asyncio, and to adopt Pipecat only if its
interruption handling saves real work.
